\centerline{\bf \Huge Текстовые задачи остаток F}

\centerline{\bf \Huge Явный объем работы}

\problem{\textit{\textbf{1}}} В помощь садовому насосу, перекачивающему 5 литров воды за 2 минуты, подключили второй насос, перекачивающий тот же объем воды за 3 минуты. Сколько минут эти два насоса должны работать совместно, чтобы перекачать 25 литров воды?


\problem{\textit{\textbf{2}}} Петя и Митя выполняют одинаковый тест. Петя отвечает за час на 10 вопросов теста, а Митя - на 16 . Они одновременно начали отвечать на вопросы теста, и Петя закончил свой тест позже Мити на 117 минут. Сколько вопросов содержит тест?


\problem{\textit{\textbf{3}}} Заказ на 180 деталей первый рабочий выполняет на 3 часа быстрее, чем второй. Сколько деталей в час делает второй рабочий, если известно, что первый за час делает на 3 детали больше?



\problem{\textit{\textbf{4}}} На изготовление 486 деталей первый рабочий затрачивает на 9 часов меньше, чем второй рабочий на изготовление 621 детали. Известно, что первый рабочий за час делает на 4 детали больше, чем второй. Сколько деталей в час делает первый рабочий?\



\problem{\textit{\textbf{5}}} Первая труба пропускает на 4 литра воды в минуту меньше, чем вторая. Сколько литров воды в минуту пропускает первая труба, если резервуар объемом 192 литра она заполняет на 4 минуты дольше, чем вторая труба?


\problem{\textit{\textbf{6}}} Первая труба пропускает на 8 литров воды в минуту меньше, чем вторая. Сколько литров воды в минуту пропускает первая труба, если резервуар объемом 660 литров она заполняет на 11 минут дольше, чем вторая труба заполняет резервуар объемом 570 литров?


\problem{\textit{\textbf{7}}} Одна тракторная бригада должна была вспахать 240 га, а другая на 35\% больше, чем первая. Вспахивая ежедневно на 3 га меньше второй бригады, первая всё же закончила работу на 2 дня раньше, чем вторая. Сколько гектаров вспахивала каждая бригада ежедневно?


\centerline{\bf \Huge Неявный объем работы}

\problem{\textit{\textbf{8}}} Один мастер может выполнить заказ за 12 часов, а другой - за 6 часов. За сколько часов выполнят заказ оба мастера, работая вместе?


\problem{\textit{\textbf{9}}} Первый насос наполняет бак за 12 минут, второй - за 14 минут, а третий - за 1 час 24 минуты. За сколько минут наполнят бак три насоса, работая одновременно?


\problem{\textit{\textbf{10}}} Две трубы наполняют бассейн за 2 часа 42 минуты, а одна первая труба наполняет бассейн за 27 часов. За сколько часов наполняет бассейн одна вторая труба?


\problem{\textit{\textbf{11}}} Игорь и Паша красят забор за 24 часа. Паша и Володя красят этот же забор за 30 часов, а Володя и Игорь - за 40 часов. За сколько часов мальчики покрасят забор, работая втроем?


\problem{\textit{\textbf{12}}} Каждый из двух рабочих одинаковой квалификации может выполнить заказ за 16 часов. Через 2 часа после того, как один из них приступил к выполнению заказа, к нему присоединился второй рабочий, и работу над заказом они довели до конца уже вместе. Сколько часов потребовалось на выполнение всего заказа?


\problem{\textit{\textbf{13}}} Первая труба наполняет резервуар на 90 минут дольше, чем вторая. Обе трубы наполняют этот же резервуар за 24 минуты. За сколько минут наполняет этот резервуар одна вторая труба?


\problem{\textit{\textbf{14}}} Двое рабочих, работая вместе, могут выполнить работу за 2 дня. За сколько дней, работая отдельно, выполнит эту работу первый рабочий, если он за 1 день выполняет такую же часть работы, какую второй - за 2 дня?


\problem{\textit{\textbf{15}}} Оператор ЭВМ, работая вместе с учеником, обрабатывает задачу за 2 ч. 24 мин. Если оператор проработает 2 ч, а ученик 1 ч, то будет выполнено $\dfrac{2}{3}$ всей работы. Сколько времени потребуется оператору и ученику в отдельности на обработку задачи?

\hspace{3mm}

\centerline{\bf \Huge Проценты}

\problem{\textit{\textbf{16}}} Митя, Антон, Гоша и Борис учредили компанию с уставным капиталом 200000 рублей. Митя внес $14 \%$ уставного капитала, Антон - 42000 рублей, Гоша - 0,12 уставного капитала, а оставшуюся часть капитала внес Борис. Учредители договорились делить ежегодную прибыль пропорционально внесенному в уставной капитал вкладу. Какая сумма от прибыли 1000000 рублей причитается Борису? Ответ дайте в рублях.


\problem{\textit{\textbf{17}}} Из данных четырех чисел первые три относятся между собой как $\dfrac{1}{5}: \dfrac{1}{3}: \dfrac{1}{20}$, а четвертое составляет $15 \%$ второго числа. Найдите эти числа, если известно, что второе число на 8 больше суммы остальных.


\problem{\textit{\textbf{18}}} Семь рубашек дешевле куртки на $2 \%$. На сколько процентов десять рубашек дороже куртки?


\problem{\textit{\textbf{19}}} Семья состоит из мужа, жены и их дочери студентки. Если бы зарплата мужа увеличилась вчетверо, общий доход семьи вырос бы на $192 \%$. Если бы стипендия дочери уменьшилась вчетверо, общий доход семьи сократился бы на $6 \%$. Сколько процентов от общего дохода семьи составляет зарплата жены?


\problem{\textit{\textbf{20}}} Букинистический магазин продал книгу со скидкой $10 \%$ с назначенной цены и получил при этом $8 \%$ прибыли. Сколько процентов прибыли первоначально полагал получить магазин?


\problem{\textit{\textbf{21}}} В 2008 году в городском квартале проживало 40000 человек. В 2009 году, в результате строительства новых домов, число жителей выросло на $8 \%$, а в 2010 году - на $9 \%$ по сравнению с 2009 годом. Сколько человек стало проживать в квартале в 2010 году?


\problem{\textit{\textbf{22}}} Цена холодильника в магазине ежегодно уменьшается на одно и то же число процентов от предыдущей цены. Определите, на сколько процентов каждый год уменьшалась цена холодильника, если, выставленный на продажу за 20900 рублей, через два года был продан за 16929 рублей.


\problem{\textit{\textbf{23}}} В четверг акции компании подорожали на некоторое число процентов, а в пятницу подешевели на то же самое число процентов. В результате они стали стоить на $36 \%$ дешевле, чем при открытии торгов в четверг. На сколько процентов подорожали акции компании в четверг?


\problem{\textit{\textbf{24}}} За год стипендия студента увеличилась на $32 \%$. В первом полугодии стипендия увеличилась на $10 \%$. Определите, на сколько процентов увеличилась стипендия во втором полугодии?


\centerline{\bf \Huge Концентрация}

\problem{\textit{\textbf{25}}} В сосуд, содержащий 5 литров 12 -процентного водного раствора некоторого вещества, добавили 7 литров воды. Сколько процентов составляет концентрация получившегося раствора?


\problem{\textit{\textbf{26}}} Смешали некоторое количество 13-процентного раствора некоторого вещества с таким же количеством 17-процентного раствора этого вещества. Сколько процентов составляет концентрация получившегося раствора?


\problem{\textit{\textbf{27}}} Смешали 9 литров 40-процентного водного раствора некоторого вещества с 11 литрами 20-процентного водного раствора этого же вещества. Сколько процентов составляет концентрация получившегося раствора?


\problem{\textit{\textbf{28}}} Виноград содержит $90 \%$ влаги, а изюм - $5 \%$. Сколько килограммов винограда требуется для получения 36 килограммов изюма?


\problem{\textit{\textbf{29}}} Имеется два сплава. Первый содержит 5\% никеля, второй - $35 \%$ никеля. Из этих двух сплавов получили третий сплав массой 225 кг, содержащий $25 \%$ никеля. На сколько килограммов масса первого сплава меньше массы второго?


\problem{\textit{\textbf{30}}} Первый сплав содержит $5 \%$ меди, второй — $13 \%$ меди. Масса второго сплава больше массы первого на 9 кг. Из этих двух сплавов получили третий сплав, содержащий $10 \%$ меди. Найдите массу третьего сплава. Ответ дайте в килограммах.

\problem{\textit{\textbf{31}}} Смешав 55-процентный и 97-процентный растворы кислоты и добавив 10 кг чистой воды, получили 65-процентный раствор кислоты. Если бы вместо 10 кг воды добавили 10 кг 50-процентного раствора той же кислоты, то получили бы 75-процентный раствор кислоты. Сколько килограммов 55-процентного раствора использовали для получения смеси?


\problem{\textit{\textbf{32}}} Имеется два сосуда. Первый содержит 100 кг, а второй - 20 кг раствора кислоты различной концентрации. Если эти растворы смешать, то получится раствор, содержащий $67 \%$ кислоты. Если же смешать равные массы этих растворов, то получится раствор, содержащий $77 \%$ кислоты. Сколько килограммов кислоты содержится в первом сосуде?

\centerline{\bf \Huge Дополнительные}

\problem{\textit{\textbf{33}}} Найдите двузначное число, если количество единиц в нем на 4 больше количества десятков, а произведение искомого числа на сумму его цифр равно 90.

\problem{\textit{\textbf{34}}} Сумма квадратов цифр двузначного числа равна 13. Если цифры поменять местами, то получится число, которое меньше данного на 9. Найдите это число.

\problem{\textit{\textbf{35}}} Сумма цифр трехзначного числа равна 11, а сумма квадратов цифр этого числа равна 45 . Если от искомого числа отнять 198, то получится число, записанное теми же цифрами в обратном порядке. Найдите число.

\problem{\textit{\textbf{36}}} При перемножении двух положительных чисел, из которых одно на 10 больше другого, ученик допустил ошибку, уменьшив цифру десятков в произведении на 4. Для проверки ответа при делении полученного произведения на меньшее число он получил в частном 39, а в остатке 22. Найдите искомые множители.

